\documentclass[sigconf,nonacm]{acmart}
% CS527 project proposal, Fall 2026. Compile with: latexmk -pdf proposal.tex
\settopmatter{printacmref=false}
\acmConference[CS527 Fall 2026]{CS527 Project Proposal}{September 2026}{Urbana, IL}
\acmYear{2026}

\usepackage{booktabs}
\usepackage{xspace}
\newcommand{\tool}{\textsc{DemandTest}\xspace}
\newcommand{\suff}{\ensuremath{\Sigma}\xspace}

\begin{document}

\title{\tool: Demand-Driven Static Discovery of Project-Specific Knowledge for Efficient Intention-Aligned Unit Test Generation}

\author{Kaleb Guo}
\affiliation{%
  \institution{University of Illinois Urbana-Champaign}
  \city{Urbana}\state{IL}\country{USA}}
\email{TODO@illinois.edu}

\begin{abstract}
Industrial unit tests are written to validate a requirement---a \emph{validation intention}---not to maximize coverage, and writing such a test requires project-specific knowledge (constructors, factories, fixtures, mocks, assertion idioms) that is scattered across thousands of files. Repository-level code agents find this knowledge by reading files one at a time, which makes them expensive in tokens and wall-clock time; existing context-injection generators retrieve fixed context patterns and have no notion of whether the retrieved context is \emph{sufficient}. We propose \tool, which addresses ICSE~2027 Industry Challenge~\#8. Given a focal method, a Java repository, and a natural-language intention, \tool derives a typed \emph{demand set} from the focal signature and the intention, resolves every demand against a one-time static project index using \emph{construction recipes}, expands context only along edges that satisfy an unresolved demand, and stops when a decidable \emph{structural context-sufficiency predicate} holds. The LLM then receives one compact knowledge packet and generates the test in a single call; mechanical failures are repaired statically and at most one LLM repair round is allowed. The research question is whether explicit construction and assertion obligations select context and decide when to stop better than similarity ranking or fixed dependency expansion at the same context budget. We will evaluate \tool against OpenHands at three iteration budgets, a faithful IntentionTest reimplementation, and a compact static-context baseline that shares our generator and repair policy, on the 13 open-source Java projects of IntentionTest's benchmark, with open-weight models, using aligned success rate over all tasks as the primary metric and tokens, first-use and amortized time, and online files inspected as the efficiency metrics. A pilot on 150 tasks across three contrasting projects decides whether the full evaluation is run. The pipeline and the shared ledger are implemented and pass offline acceptance tests; no experimental results exist yet.
\end{abstract}

\maketitle

\section{Problem}
\label{sec:problem}

\paragraph{Background.}
Coverage is not what developers write tests for. In practice a test validates that a requirement, a feature scenario, or an expected behavior has been implemented; Qi et al.~\cite{qi2025intentiontest} call this a \emph{validation intention} and instantiate it as a short semi-structured description: an \emph{objective} ($\le$50 words), optional \emph{preconditions}, and optional \emph{expected results}. Turning such an intention into an executable unit test is costly in industry because the required project-specific knowledge---which constructor or factory produces a valid input, which fixtures and mocks the project uses, which helper asserts the expected state---is spread over a large codebase and is rarely textually similar to the intention. ICSE~2027 Industry Challenge~\#8~\cite{icse27challenge8} states the problem precisely: \emph{given a focal method, a large industrial codebase, and a natural-language validation intention, how can we efficiently generate an executable, project-specific unit test that faithfully validates the intended behavior?} The challenge lists three obstacles---efficient discovery of testing knowledge, assessing whether the retrieved knowledge is \emph{sufficient}, and balancing effectiveness, efficiency, and confidentiality---and five requirements: intention alignment (R1), scalability to repositories with more than 10{,}000 source files (R2), at least 50\% fewer tokens and 20\% less wall-clock time than strong repository-agent baselines at comparable quality (R3), deployability with local or private LLMs (R4), and evaluation on compilation rate, execution rate, semantic alignment, developer acceptance, token cost, response time, and number of inspected files (R5).

\paragraph{Why existing approaches fall short.}
Search-based tools such as EvoSuite~\cite{fraser2011evosuite} maximize structural coverage and cannot target an intention. Prompt-only LLM generators such as ChatTester~\cite{yuan2024chattester} hallucinate project APIs and miss setup logic. A recent line of work injects project context obtained by static analysis or a language server---KTester~\cite{li2026ktester}, RATester~\cite{yin2025ratester}, APT~\cite{zhang2024apt}, CATGen~\cite{chen2026catgen}---and IntentionTest~\cite{qi2025intentiontest} retrieves a similar existing test and edits it under the intention with up to four LLM refinement rounds. These systems raise compilation and pass rates, but they either look up context reactively inside a multi-turn loop (RATester), or inject fixed context patterns with no criterion for when enough has been collected, and none of them is measured against repository agents on the challenge's efficiency metrics; IntentionTest reports no token or latency numbers at all. TestTailor~\cite{zhou2026testtailor}, a coverage-driven generator, shows that reusing the existing test closest to the target and telling the model where it diverges is the single largest lever in its ablation, three times the effect of symbolic path constraints; it generates no assertions and names project-specific object construction as its open problem, which is exactly what our recipes provide. Within its 2021--2025 window, a 2026 systematic review of 116 studies~\cite{zhang2026slr} lists IntentionTest and the ProjectTest benchmark (since renamed MultiFileTest~\cite{wang2026multifiletest}) as the only project-level systems and inventories no efficiency metric; the exceptions we know of appeared after that window: CATGen~\cite{chen2026catgen} reports tokens and generation time against other LLM-based generators, and Ferrer and Chicano~\cite{ferrer2026context} measure the cost of dependency-retrieval strategies on industrial code, neither against a repository agent. The review's identified needs, program-analysis-driven context selection that fits the model's window and dependency-aware prompts that state how to obtain values, are what \tool computes statically. Our claim is therefore not ``static analysis plus compact context plus fewer calls'', which KTester, CATGen, and Ferrer and Chicano share, but that typed obligations decide \emph{which} context to retrieve and \emph{when to stop}, and that this beats similarity ranking and fixed expansion at the same budget. Repository agents such as SWE-agent~\cite{yang2024sweagent} and OpenHands~\cite{wang2025openhands} are strong on repository-level test generation~\cite{mundler2024swtbench} but pay for it with long tool-call trajectories, and cost is now a first-class concern in agent design~\cite{wang2025efficientagents}. Agentless~\cite{xia2025agentless} showed for program repair that a fixed localize--generate--validate pipeline can rival agents at a fraction of the cost. We ask the analogous question for intention-driven test generation, with one addition the challenge asks for explicitly: a decidable answer to ``is the collected knowledge sufficient?''

\paragraph{Problem statement.}
Design a generator that, for a focal method $m$ in class $C$ of a Java repository $R$ and an intention $I$, produces a JUnit test that compiles, executes, and validates $I$, while minimizing LLM tokens, LLM calls, wall-clock time, and the number of repository files inspected, and that runs with open-weight models. We do not study program repair, vulnerability detection, coverage-driven generation, or general-purpose agents, and we never claim to find real bugs: correctness is measured against the real code and against held-out human-written tests.

\section{Solution}
\label{sec:solution}

\paragraph{Thesis.}
Most of what an agent spends tokens on is \emph{finding} testing knowledge. That knowledge is largely determined by types: to call $m$ one needs a receiver of type $C$ and arguments of the parameter types; to set up $m$ one needs the collaborators $m$ reads; to assert one needs an observable of the return or state type. These needs can be computed before any LLM call, resolved against a static index of the project's construction paths and test idioms, and checked for sufficiency without an LLM. \tool is that computation followed by one generation call. Table~\ref{tab:stages} summarizes the pipeline; the full specification, including the index schema, module contracts, prompts, and acceptance tests, is in the repository (\url{https://github.com/TODO/demandtest}, \texttt{docs/SPEC.md}).

\begin{table}[t]
  \caption{\tool stages and the LLM calls each may make.}
  \label{tab:stages}
  \small
  \begin{tabular}{@{}lp{0.66\linewidth}c@{}}
    \toprule
    Stage & What it does & LLM calls \\
    \midrule
    S0 & Project index, once per repository & 0 \\
    S1 & Demand set $D(m,I)$ from signature and intention & 0 \\
    S2 & Bounded expansion until \suff holds & 0 \\
    S3 & Generation from one knowledge packet & 1 \\
    S4 & Static repair; compile and run & 0 \\
    S5 & Semantic repair, only on failure & $\le 1$ \\
    \bottomrule
  \end{tabular}
\end{table}

\paragraph{S0: Project index.}
A one-time pass with JavaParser and its symbol solver~\cite{javaparser} records, for every type, its constructors, static factories, builders, singletons (including enum constants), public observables (getters, \texttt{equals}, \texttt{toString}, public fields), the fields each method reads and writes, and, for every existing test, its fixtures, helper methods, resolved callees, and the assertion and mocking libraries it uses. Every entity carries the file it lives in, so file accounting is exact.

\paragraph{S1: Demand set.}
$D(m,I)$ contains a \emph{receiver} need for $C$ (unless $m$ is static), one \emph{argument} need per parameter, annotated with the sentence of the preconditions that mentions the parameter, one \emph{setup} need per field of $C$ that $m$ reads, \emph{oracle} needs derived from the expected results (an exception need if they mention throwing, otherwise a return-value need and, if they describe state change or $m$ is void, a state need; an interaction need if they mention delegation), and an \emph{idiom} need for the project's test framework, assertion library, and mocking library. No LLM is involved.

\paragraph{S2: Recipes, sufficiency, bounded expansion.}
A \emph{recipe} is a way to obtain a value of a type: a JDK literal, a fixture or helper from an existing test, a constructor, a factory, a builder, a singleton, or a mock (only for interfaces and abstract classes when the project already mocks). Its cost is the number of parameters that are not literal-resolvable, summed over the sub-recipes chosen for those parameters; ties prefer how the project already does it (fixture before constructor before factory before mock). Resolution is demand-driven in the sense of demand-driven program analysis~\cite{heintze2001demand}: only the types a need asks for are ever resolved, recursively to a bounded depth, with memoization. The \emph{structural context-sufficiency predicate} \suff$(D,\mathit{ctx})$ holds when every receiver, argument, and setup need has a recipe whose entities are in the context $\mathit{ctx}$ and none of whose leaves is unresolved, every oracle need has an observable in $\mathit{ctx}$ (or a mock is admissible), and the idiom is known. \suff is structural: it establishes that every value is obtainable and that an observable exists, not that the context suffices for a semantically aligned test. It does not know the relation between two arguments an intention constrains jointly, the receiver state a scenario requires, or which observable distinguishes the intended outcome from an unrelated failure. S1 records such cases as \emph{semantic gaps} (relational and state sentences rendered verbatim but unbound, sentences it cannot bind, several arguments of one type); a zero count means no \emph{detected} gap. Because the token budget may cut the rendered packet, \suff is claimed only for evidence the model actually receives: each run's status is computed on the rendered packet as \emph{sufficient}, \emph{sufficient-with-gaps}, \emph{budget-limited}, or \emph{fallback}, and every quality metric is reported by status. Every unresolvable need carries a reason (no construction path, dependency injection, reflection, resource files, off-index factories, erased generics) whose distribution is reported, since JavaParser with its symbol solver is a starting point and not a complete resolver. Expansion starts from the focal signature and, while \suff fails and a file budget $B_f$ (default 12) and a token budget $B_t$ (default 2{,}000) are not exhausted, picks an unresolved need, resolves it over the whole index, and adds the recipe's entities and files to $\mathit{ctx}$. If \suff still fails, \tool falls back to the two existing tests whose already-satisfied needs overlap most with $D$, a static analog of TestTailor's path-proximal tests~\cite{zhou2026testtailor}, shows each with its \emph{demand diff} (what it already obtains, what it does not assert), and marks the run as \emph{fallback}, so that the predicate's accuracy can be measured. Oracle needs also carry a syntactic trigger hint mined from the focal body (the guards enclosing the throw or return the intention asks for), a static, single-pass counterpart of the trace-derived guard conditions exLong~\cite{zhang2025exlong} extracts for exceptional-behavior tests. Every step is logged as a trace.

\paragraph{S3--S5: One call, static repair, one repair call.}
The context is rendered into a deterministic \emph{knowledge packet}: intention, focal source, one line per resolved need with the recipe as a code fragment and the signatures it depends on, observables for the oracle, project facts, and one assertion-style example. A single temperature-0 call produces the test class; a parse failure counts as a failed run, with no retry. Static repair then fixes the package declaration (which also grants package-private access), missing imports resolved through the index, missing \texttt{throws} clauses, and JUnit-4/5 idiom mismatches---following the observation of CATGen~\cite{chen2026catgen} that program-analysis post-processing removes most iterative LLM repair; we adopt this step and do not claim it. Only if the test compiles but fails, or javac diagnostics survive, is one repair call made with the trimmed diagnostics and the packet, opening no new files.

\paragraph{Harness and deployability.}
All systems, including baselines, write to one SQLite ledger (\texttt{llm\_calls}, \texttt{file\_access}, \texttt{results}); every number reported is a query over it. The LLM client is an OpenAI-compatible HTTP call with no other dependency, so open-weight models on a private endpoint satisfy R4. Runs are idempotent per task and configuration and resume after interruption.

\paragraph{What is new.}
Relative to the closest work, \tool contributes (i)~a demand representation, typed construction and assertion obligations derived from the signature and the intention before any LLM call, that decides which context to retrieve; (ii)~a decidable, explicitly structural stopping rule with bounded expansion, evaluated on the same tasks against continuing past the rule and against fixed budgets; (iii)~evidence, at equal context budget and with the same generator and repair policy, on whether obligations select context better than a compact static context, similarity ranking, or fixed dependency expansion; and (iv)~a comparison against a repository agent at several budgets on the challenge's efficiency metrics, with index and online costs separated. Static analysis, compact context, and few calls are shared with KTester, CATGen, and Ferrer and Chicano and are not claimed.

\section{Evaluation Setup}
\label{sec:eval}

\paragraph{Subjects.}
The 13 open-source Java projects of IntentionTest~\cite{qi2025intentiontest} (51--1{,}099 tests each), for comparability, plus one or two active repositories with more than 10{,}000 source files (candidates: Elasticsearch, Apache Flink, Spring Framework), chosen after verifying build feasibility, to address R2. The Java projects of MultiFileTest~\cite{wang2026multifiletest} are a candidate published project-level set, to be added only if they do not overlap with IntentionTest's subjects. No industrial repository is available to us. The nearest proxies, reported as proxies, are (a)~at least 30 intentions written independently by two people from the focal method and its class only, never from a test, and (b)~a blinded rating of generated tests by three to five developers on acceptance, edits required, and usefulness. Both are a stated threat to R5's industrial evaluation, not a substitute for it.

\paragraph{Tasks and ground truth.}
Following IntentionTest's benchmark construction, an intention is reverse-engineered by an LLM from each existing test that exercises exactly one focal method (objective $\le$50 words, preconditions plus expected results $\le$200 words, no more than five code tokens); 10\% are checked by two people with agreement reported. That sample validates the intentions, not the output judge, which has its own sample. Trivial focal methods (getters, setters, $\le$3 LOC) are dropped. Leakage control: no information derived from the held-out test enters any system's index or context. Its entry, fixtures, callees, and source are hidden from every query, and the index holds no precomputed summaries or retrieval features that could survive exclusion. Class-level fixtures shared with sibling tests remain, as legitimate neighboring context, and every quality metric is reported separately for tasks with a close visible clone (a sibling test calling the same method, or demand overlap above 0.8) and tasks without one. The agent's checkout has the reference method removed and its git history stripped.

\paragraph{Baselines and ablations.}
The primary baseline, required by R3, is OpenHands~\cite{wang2025openhands} in headless mode with the same model and the intention and focal location as the task, at three iteration budgets (10, 30, 60) after checking on five tasks that the chosen open-weight model drives its tools competently; its trajectories are parsed into the same ledger. A single cap would not establish a fair comparison. Secondary baselines are a \emph{simplified} reimplementation of IntentionTest's retrieve-and-edit pipeline on our harness, labeled as such because it omits their crucial-fact discrimination stage, and a \emph{compact static-context} baseline: the focal method plus the signatures of directly referenced types and the class's fields under the same how-to-obtain heading, cut to the same allowance, with the same generator and repair policy. That whole-system comparison does not isolate selection, so two controlled experiments do. The \emph{selection} experiment holds the candidate pool, rendering, auxiliary hints, stopping policy, and generator fixed at several context allowances and changes only the selector: typed obligations, similarity ranking, or dependency distance. The \emph{stopping} experiment holds the selector fixed and varies the stopping policy over genuinely nested packets produced by a deterministic continuation rule, with the four-entity extension preselected as the comparator; the full focal class is a separate representation ablation, since it is not a superset of a packet holding external factories and fixtures. We also report intention sensitivity, the fraction of tasks where changing the intention changes the selected entities, because recipe resolution is type-driven and only cue extraction, annotation, and related-test ranking depend on the intention. Further ablations: the recipe tie-break swapped and the return-value trigger hint removed. SWE-agent~\cite{yang2024sweagent} time permitting.

\paragraph{Models.}
Main experiments use open-weight models (Qwen3-Coder or DeepSeek-V3 class) through an OpenAI-compatible endpoint, which is what R4 requires; a closed model is run on a subset as an upper-bound reference.

\paragraph{Metrics (R5).}
The primary metric is \emph{aligned success rate} over all attempted tasks: the fraction whose test compiles, passes, executes the focal method, and receives the top alignment score. Focal execution is measured dynamically with JaCoCo on the generated test, because a static call reference may be unreachable or mocked; the static \emph{target-hit} check (the test calls $m$ and contains the oracle the intention asks for), after TestTailor's coverage-accuracy metric~\cite{zhou2026testtailor}, is reported next to it as the cheap proxy. Alignment is an LLM judge on a 0--2 rubric whose agreement with blinded human ratings of 100 generated tests, drawn across systems and projects with two raters, is reported before the judge's numbers are used; we keep it separate from pass rate because LLM-written oracles tend to encode the behaviour a program has rather than the behaviour it should have~\cite{konstantinou2024oracles}. Aligned-given-pass is a diagnostic only. Also reported: compilation and pass rates; mutation score with PIT~\cite{coles2016pit} restricted to the focal method's lines on a sampled subset; usable tests per LLM call; whether a repair round weakened the test (fewer assertions, lost oracle, or lost target hit); prompt and completion tokens as reported by the provider, never estimated; number of LLM calls; and cost per aligned test. Efficiency accounting is split into four buckets with the same boundaries for every system and with failed attempts included in every total: index cost (S0 wall-clock, memory, files scanned, once per repository), online retrieval (files whose entities enter the packet, and packet tokens; for the agent, every file a tool call reads or greps into), generation and repair, and execution. Time per task is reported both first-use, indexing plus online, and amortized, indexing divided by the number of tasks on the repository plus online.

\paragraph{Research questions.}
RQ1: Is aligned success non-inferior to the agent's? Comparable quality is predefined as non-inferiority: the lower 95\% confidence bound on the paired difference against OpenHands at 30 iterations must exceed $-0.05$, with sensitivity reported at $-0.03$ and $-0.02$. Five points is an engineering tradeoff, stated as such: the intended use is a first-draft generator that a developer reviews, and the savings R3 asks for are worth five reviewed drafts per hundred that need more work. A non-significant difference is not parity. RQ2: Do tokens, first-use and amortized time, and online files drop past the R3 thresholds at every agent budget? RQ3: Under a common context allowance, with the same pool, rendering, and generator, does the obligation selector beat similarity ranking and dependency distance? RQ4: How sensitive are results to model size, to objective-only intentions, to independently written intentions, and to the clone-available versus clone-free strata? RQ5: Does the stopping rule stop at the right point? On the same tasks, stopping at \suff is compared with the nested extensions by two and four entities and with filling the allowance, and a manual analysis of 90 failed tasks drawn from every status classifies each as premature stop, intention-binding failure, generation failure with adequate context, or repair failure. Comparing sufficient against fallback tasks alone would not establish predicate accuracy, since the former may simply be easier.

\section{Result Analysis}
\label{sec:results}

There are no experimental results yet; this section reports what is in place and how the results will be read.

\paragraph{Current status.}
The design is fixed in a specification (\texttt{docs/SPEC.md}) that defines the demand set, recipes, the structural predicate and its boundary, the expansion algorithm, the packet format, the fixed prompts, the index schema, the ledger schema, the CLI, the task-construction and leakage-control procedure, the baseline protocol, and per-module acceptance tests. The indexer, the S1--S5 pipeline, the ledger, and the accounting client are implemented and pass 88 offline acceptance tests; cron-utils is indexed. No task has yet been run against a model. The pilot protocol (\texttt{docs/PILOT.md}) fixes 150 tasks across three contrasting projects, two controlled experiments for selection and stopping, a 45-task paired agent comparison, the manual analyses, and a three-way decision per question: proceed, redesign, or inconclusive. Its statistics are paired at task level, tasks sharing a focal method are one unit, the three projects are fixed case studies reported separately without a cluster bootstrap over three clusters, and generalization across projects is claimed only in the full evaluation, whose sample is planned from the pilot's observed paired disagreement rate.

\paragraph{How results will be read.}
The headline artifact is a Pareto plot of aligned success rate against tokens per task with confidence intervals, one point per \tool configuration and per agent budget, and a table of the R5 metrics per system with index and online costs separated. The R3 thresholds are necessary, not sufficient; the pilot's three-way decisions are the criterion for running the full evaluation, and an inconclusive interval leads to sample planning or a narrower claim, not to redesign. Online files inspected is expected to show the largest gap, since after indexing \tool touches a file only when a need requires it; the index scan is reported as first-use cost. Failure modes we expect and will report: \suff holds while the call fails, which RQ5 separates into premature stops and generation failures; intentions with relational or state preconditions that S1 cannot bind; setup over-approximation on large classes; and large repositories that do not build.

\section{Future Plan}
\label{sec:plan}

The course project concludes in mid-November, aligned with the challenge's solution-paper deadline (November~13, 2026). The next month is the critical path.

\paragraph{Through October~5.}
Model endpoint; ten hand-written cron-utils tasks run end to end and the first real failures fixed; the intention script; dynamic focal execution; the stopping-rule and ablation flags; the compact static-context baseline; OpenHands on five tasks.

\paragraph{October~6--19: the pilot.}
150 tasks across three contrasting projects, 30 of them with independently written intentions; all \tool conditions and the compact static-context baseline; OpenHands at three budgets and the IntentionTest reimplementation on a 45-task subset; the missing-fact and judge-validation analyses; the proceed / redesign / inconclusive decision on whether obligations select better context, whether stopping preserves aligned success, and whether savings survive honest accounting.

\paragraph{October~20 -- November~13.}
If go: scale to the 13 projects and one large repository, mutation testing, the variance protocol, analysis, figures, and the write-up. If redesign: fix the failing component and report the pilot as the result; if inconclusive: plan the sample or narrow the claim. Collaborators from the lab take the indexer, the harness, and the baseline-and-evaluation tracks respectively; a weekly sync reviews the ledger.

\bibliographystyle{ACM-Reference-Format}
\bibliography{proposal}

\end{document}
